% ==========================================================
% PALETA DE CORES
% ==========================================================

\definecolor{AzulPrincipal}{HTML}{173F5F}
\definecolor{AzulClaro}{HTML}{20639B}
\definecolor{Turquesa}{HTML}{3CAEA3}
\definecolor{Verde}{HTML}{2A9D8F}
\definecolor{Amarelo}{HTML}{F6C85F}
\definecolor{Laranja}{HTML}{F4A261}
\definecolor{Vermelho}{HTML}{E76F51}
\definecolor{Roxo}{HTML}{6C63A8}
\definecolor{VerdeDestaque}{HTML}{3A9D5D}
\definecolor{VerdeDestaqueFundo}{HTML}{F0FAF2}

\definecolor{Fundo}{HTML}{F7F9FC}
\definecolor{CinzaTexto}{HTML}{4A5568}
\definecolor{CinzaClaro}{HTML}{E9EEF4}
\definecolor{VermelhoSolucao}{HTML}{C94C4C}
\definecolor{VermelhoSolucaoFundo}{HTML}{FDF2F2}
\definecolor{TurquesaTeorema}{HTML}{159A9C}
\definecolor{TurquesaTeoremaFundo}{HTML}{F0FAFA}

% ==========================================================
% CAIXA DE QUESTÃO
% ==========================================================

\newtcolorbox{caixaquestao}{
    enhanced,
    breakable,
    colback=white,
    colframe=CinzaClaro,
    boxrule=0.7pt,
    arc=8pt,
    left=12pt,
    right=12pt,
    top=12pt,
    bottom=12pt,
    before skip=12pt,
    after skip=18pt,
    borderline west={4pt}{0pt}{AzulPrincipal}
}

% ==========================================================
% CAIXA DE DESAFIO
% ==========================================================

\newtcolorbox{caixadesafio}{
    enhanced,
    breakable,
    colback=Amarelo!12,
    colframe=Amarelo!70!black,
    boxrule=0.8pt,
    arc=8pt,
    left=14pt,
    right=14pt,
    top=12pt,
    bottom=12pt,
    title={\sffamily\bfseries DESAFIO},
    coltitle=AzulPrincipal,
    fonttitle=\bfseries
}

% ==========================================================
% CAIXA DE DICA
% ==========================================================

\newtcolorbox{caixadica}{
    enhanced,
    breakable,
    colback=Turquesa!8,
    colframe=Turquesa,
    boxrule=0.7pt,
    arc=8pt,
    left=14pt,
    right=14pt,
    top=10pt,
    bottom=10pt,
    title={\sffamily\bfseries DICA},
    coltitle=AzulPrincipal,
    fonttitle=\bfseries
}

% ==========================================================
% CAIXA DE INFORMAÇÃO
% ==========================================================

\newtcolorbox{caixainfo}{
    enhanced,
    breakable,
    colback=AzulPrincipal!5,
    colframe=AzulPrincipal,
    boxrule=0.7pt,
    arc=8pt,
    left=14pt,
    right=14pt,
    top=10pt,
    bottom=10pt
}

% ==========================================================
% COMANDO PARA NÍVEL
% ==========================================================

\newcommand{\nivel}[1]{
    \textcolor{CinzaTexto}{\small\textbf{Nível: }}
    \foreach \n in {1,...,3}{
        \ifnum\n<#1+1
            \textcolor{Laranja}{\Large$\bullet$}
        \else
            \textcolor{CinzaClaro}{\Large$\bullet$}
        \fi
    }
}

% ==========================================================
% IDENTIFICAÇÃO DA QUESTÃO
% ==========================================================

\newcommand{\cabecalhoquestao}[2]{
    \begin{tikzpicture}
        \node[
            circle,
            fill=AzulPrincipal,
            text=white,
            minimum size=0.85cm,
            inner sep=0pt,
            font=\bfseries
        ] {#1};

        \node[
            anchor=west,
            font=\small\sffamily\bfseries,
            text=AzulPrincipal
        ] at (0.6cm,0) {QUESTÃO};

        \node[
            anchor=west,
            font=\small\sffamily,
            text=CinzaTexto
        ] at (2.0cm,0) {#2};
    \end{tikzpicture}
}

% ==========================================================
% QUESTÃO
% ==========================================================

\newcounter{questao}

\newcommand{\questao}[3]{
    \stepcounter{questao}

    \begin{caixaquestao}

        \cabecalhoquestao{\thequestao}{#1}

        \vspace{8pt}

        #2

        \vspace{8pt}

        #3

    \end{caixaquestao}
}

% ==========================================================
% ALTERNATIVAS
% ==========================================================

\newenvironment{alternativas}{
    \begin{enumerate}[
        label=\textbf{\Alph*)},
        labelsep=0.5em,
        itemsep=5pt
    ]
}{
    \end{enumerate}
}

% ==========================================================
% IDENTIFICAÇÃO PEDAGÓGICA
% ==========================================================

\newcommand{\metadados}[3]{
    \vspace{5pt}

    \begin{tikzpicture}
        \node[
            fill=AzulPrincipal!7,
            rounded corners=5pt,
            inner xsep=8pt,
            inner ysep=4pt,
            text=AzulPrincipal,
            font=\scriptsize\sffamily
        ] {
            \textbf{Conteúdo:} #1
            \hspace{10pt}
            \textbf{Habilidade:} #2
            \hspace{10pt}
            \textbf{#3}
        };
    \end{tikzpicture}
}

% ==============================================================
% CAIXA SOLUÇÃO
% ==============================================================

\newtcolorbox{caixasolucao}{
    colback=VermelhoSolucaoFundo,
    colframe=VermelhoSolucao,
    boxrule=0pt,
    leftrule=3pt,
    arc=2pt,
    left=8pt,
    right=8pt,
    top=8pt,
    bottom=8pt,
    enhanced,
    breakable,
}

% ==================================================================
% CAIXA TEOREMA
% ==================================================================

\newtcolorbox{Teorema}[1]{
    colback=TurquesaTeoremaFundo,
    colframe=TurquesaTeorema,
    boxrule=0pt,
    leftrule=4pt,
    arc=2pt,
    left=10pt,
    right=10pt,
    top=8pt,
    bottom=8pt,
    enhanced,
    before upper={%
        \textcolor{TurquesaTeorema}{\textbf{Teorema — #1}}
        \par\smallskip
    },
}